% =====================================================================
% No-Go Theorem (canonical statement)
% =====================================================================

\chapter{No-Go Theorem (Canonical)}
\label{ch:nogo-canonical}

\section{Definition}

\begin{theorem}[No-Go]
\label{thm:nogo-canonical}
Assume \textbf{only} the following data:
\begin{enumerate}[label=(\roman*)]
  \item the residue-class structure of the ternary map \(T_3\) (including the
        published charge-preserving correction that restores
        \(Q=n\bmod 9\));
  \item the topology of the \(243\) Kaluza--Klein towers;
  \item the democratic partition of the flux integer \(4880\) that produces
        the initial seeds \(20\) and \(21\);
  \item any completion of the map on the previously impossible residue classes
        that is fixed solely by minimality of the charge defect among
        \(k\in\{0,1,2\}\) (the completed map \(T^\sharp\)).
\end{enumerate}
Do \textbf{not} assume any numerical bound that already contains the integer
\(539\), the puncture count \(61\), or the period \(539.9\).

Then the following hold:
\begin{enumerate}[label=(\alph*)]
  \item Flux democracy cannot break the circularity of the contraction-factor
        argument that sets \(\lambda=\ln 3/539\).
  \item The specific Banach rate \(\lambda=\ln 3/539\) cannot be derived from
        the data above.
  \item No unique generational dictionary that inserts the number \(539\)
        (in particular the windings \(w_j=539+61j\)) is forced by those data.
\end{enumerate}
Equivalently, the following programme claims are \textbf{blocked} (ruled out as
consequences of the assumed data alone):
\begin{enumerate}[label=(\alph*$'$)]
  \item Flux democracy breaks the circularity of \(\lambda=\ln 3/539\).
  \item \(\lambda=\ln 3/539\) follows from the assumed data.
  \item The dictionary \(w_j=539+61j\) is forced by the assumed data.
\end{enumerate}
\end{theorem}

\section{Derivation and equations}

The local logarithmic expansion factors of the ordinary map are
\[
\chi_0=\ln\frac13,\qquad
\chi_1=\ln\frac43,\qquad
\chi_2=\ln\frac23.
\]
Their uniform average is the a priori negative number
\begin{equation}
\frac13(\chi_0+\chi_1+\chi_2)
=\frac13\ln\frac{8}{27}
\approx -0.405.
\label{eq:unres-avg}
\end{equation}
After the charge-preserving completion \(T^\sharp\) the leading ratio on every
branch-2 step tends to \(1/3\). Under the natural \(3\)-adic equidistribution
consistent with the \(243\)-tower average the stationary expectation becomes
\begin{equation}
\mathbb{E}_\pi[\chi]
=\frac23\ln\frac13+\frac13\ln\frac43
=\ln\Bigl(\frac{4^{1/3}}{3}\Bigr)
\approx -0.6365
<0.
\label{eq:Epi-nogo}
\end{equation}
The associated mean contraction rate is
\begin{equation}
\lambda_{\mathrm{mean}}
=\frac{4^{1/3}}{3}
\approx 0.529
<1.
\label{eq:lmean}
\end{equation}
The non-circular e-fold bound that follows from this rate and the flux integer is
\begin{equation}
N_\star
=\Bigl\lceil
\frac{\ln 4880}{\bigl|\mathbb{E}_\pi[\chi]\bigr|}
\Bigr\rceil
=14.
\label{eq:Nstar-nogo}
\end{equation}
None of these quantities equals \(539\), and none of the steps that produce them
inserts the number \(539\).

\section{Proof}

\begin{proof}[Proof of Theorem~\ref{thm:nogo-canonical}]
(1)~The democratic partition fixes only the initial seeds. It supplies no
information about the mean size of the charge-correcting exponent along later
iterates that would determine a Lipschitz constant already containing \(539\).
Therefore democracy cannot break the circularity of the argument that sets
\(\lambda=\ln 3/539\). This is~(a).

(2)~The residue Markov structure together with the completed map \(T^\sharp\)
yields a strictly negative stationary mean \(\mathbb{E}_\pi[\chi]\approx -0.6365\)
as in \eqref{eq:Epi-nogo}. The derivation uses only residue probabilities and
the asymptotic ratio \(1/3\); it does not use the integer \(539\). The resulting
e-fold depth is \(N_\star=14\) as in \eqref{eq:Nstar-nogo}, not \(539\). The rate
actually obtained is \(\lambda_{\mathrm{mean}}\) in \eqref{eq:lmean}, not
\(\ln 3/539\). This is~(b).

(3)~Suppose, for contradiction, that the same data forced the identification
\(\lambda=\ln 3/539\). Then the Banach fixed-point estimate would recover a step
count of order \(539\). But the only Lipschitz-scale constant that the data
actually determine is \(\lambda_{\mathrm{mean}}\approx 0.529\), whose Banach /
e-fold bound is of order \(14\). The identification \(\lambda=\ln 3/539\) is
therefore an additional assumption, not a consequence. Any generational
dictionary that embeds the number \(539\) (in particular \(w_j=539+61j\))
inherits the same extra assumption and is likewise unforced. This is~(c).

(4)~The long resonant trajectory of length \(539\) is maintained only by
constraints that lie outside the listed data: a fixed iteration count, a
holographic window, and phase-locking to a pre-chosen period. Those constraints
are not implied by residue structure, tower topology, or flux democracy.

Hence claims (a)--(c) hold. The no-go stands.
\end{proof}

\section{Corollary}

\begin{corollary}[Empirical search remains permissible]
\label{cor:emp-nogo}
The empirical search for a resonant period near \(539.9\) remains permissible
precisely because it does not claim to derive that period from the data covered
by Theorem~\ref{thm:nogo-canonical}. The period is treated as a hypothesis to be
tested after a non-circular spectral estimate has been obtained
(periodogram / multitaper / Lomb--Scargle; bootstrap free of \(539.9\);
compatibility only afterward).
\end{corollary}

\begin{remark}[Depth macro split]
\label{rem:depth-split-nogo}
Throughout this book the ACE e-fold depth is \(N_\star=14\) and the model HQCC
depth is \(\sigma=539\) (equivalently \(N_{\mathrm{HQCC}}\)). These symbols are
never identified. See Chapter~\ref{ch:provenance}.
\end{remark}

% end
